For $x>0$ let $W(x)$ be the least waste of a packing of unit squares, with pairwise disjoint interiors and rotations allowed, in a square of side $x$. Erd\H{o}s and Graham \cite{EG75} showed $W(x)=O(x^{7/11})$. The best exponent is now $3/5$ \cite{Bui25,McC26} (the earlier claim \cite{CG20} contained an error \cite{AB26}). From below, Roth and Vaughan \cite{RV78} proved that $W(x)$ is not $o(x^{1/2})$. For sides just below an integer $k$, let $s(N)$ be the side of the smallest square containing $N$ unit squares and $\cs(k)$ the largest $c$ with $s(k^2-c)=k$; Ryu \cite{Ryu26} proved $\cs(k)=O(k^{3/8})$.

These constructions cut the square into an aligned grid, long thin strips, and trapezoids whose two parallel sides differ by much less than their height, and they pack each piece on its own. The exponents $3/5$ and $3/8$ come from balancing the strips against the trapezoids. Bui conjectured \cite[Conjecture~1]{Bui25} that a right trapezoid of height $x$ whose bases differ by $x^{1-\gamma}$, $0<\gamma<1$, cannot be packed with waste $o(x^{1-\gamma/2})$.

This note proves lower bounds for such pieces taken on their own. Theorem~\ref{thm:B} is Bui's conjecture with an explicit constant, at any width; his height $x$ is our $L$ and the difference of his bases is our $R$. Theorem~\ref{thm:S} is a different inequality, for regions under any concave ceiling, and it is useful when the region is thin. For a sloped ceiling it gives $L/\sqrt M$ (Corollary~\ref{cor:slope}), which is larger than $\sqrt{LR}$ exactly when $M<L/R$. For rectangles it gives the order of the Roth--Vaughan bound at half-integer sides with the constant $5\cdot10^{-4}$ (Corollary~\ref{cor:rect}).

Neither theorem is a lower bound for $W(x)$ or for $\cs(k)$. They bound regions packed on their own, and a packing of a square need not split into such regions. What they say about the constructions is in Section~\ref{sec:consequences}: given the decompositions described in \cite{Bui25,McC26,Ryu26}, the trapezoid fillers there are optimal up to constants, and those frames cannot give exponents below $3/5$ and $3/8$.

\subsection*{Method and credits}
Both proofs start a path at each point of the floor and follow it upward, crossing each square along its normal. A path that reaches the ceiling satisfies a height identity (\eqref{eq:vheight} and \eqref{eq:sheight}): the height of the ceiling where it arrives is the number of squares crossed, minus a deficit $\sum(1-\cos\varphi_i)$ coming from their tilts, plus the length spent in gaps. Paths that do not arrive are charged to the waste.

For Theorem~\ref{thm:S} a path moves vertically in a gap. The flow is then injective and preserves measure. Paths that run into the lateral side of a square are bounded by a Cauchy--Schwarz argument (Lemma~\ref{lem:vlateral}), and that term is what limits the theorem to thin regions.

For Theorem~\ref{thm:B} a path keeps, in a gap, the direction of the square it has just left. Entries through lateral sides, and collisions between the gaps above two squares, then cost only a tilt difference (Lemmas~\ref{lem:slateral} and~\ref{lem:conflicts}). Under a sloped ceiling the gap term of the arriving paths grows along the floor, so across the boundaries between neighbouring families of arriving paths it must drop by about $R$ in total (Lemma~\ref{lem:drops}). Each drop is paid for by a tilt difference times a chord length, or else forces a long step along the floor (Lemma~\ref{lem:skip}), and the sum of those products is bounded by the waste through horizontal chords (Lemma~\ref{lem:sideways}).

The waste lemma for two squares of different tilt (Lemma~\ref{lem:fundamental}) is due to Roth and Vaughan \cite{RV78}. The proof given here follows Ryu \cite[Proposition~3.3]{Ryu26b}, adapted to packings whose squares may touch. Following paths through the squares along their normals, the height identity, and the rule of keeping direction in gaps are Ryu's, from his lower bound of order $k^{1/4}$ for $\cs(k)$ \cite{Ryu26b}. Proposition~3 of \cite{Bui25} is an earlier lower bound of this kind, for a thin parallelogram. We could not consult \cite{RV78} directly and know it through secondary accounts \cite{Dan26c}; nothing below depends on it, since Lemma~\ref{lem:fundamental} is proved in full. The question came from Evan Daniel's notes on the exponent of $\cs(k)$ \cite{Dan26c}, which ask for a local lower bound of this kind.

\subsection*{Status}
The proofs in this note were found and written by Claude (Anthropic), working as a research agent under the author's direction, in one day. They have been reviewed only by language models, six blind passes in all, and tested by floating-point simulation. Section~\ref{sec:checks} gives the details. The author has not verified every step by hand, no other person has read the proofs, and nothing is formalised. The statements should be read as claims.

\subsection*{Setting}
